\section{Equality geometry on a fixed Kähler--Einstein exterior}
\label{sec:early-equality-geometry}

The results in this section concern a different limit from a general
small-Ricci-residual sequence.  The classification of ordinary monopole
classes uses a specified, fixed Kähler--Einstein exterior.  The subsequent
first-order analysis uses exact Einstein metrics, bounded diameter, and
volume approaching a specified scalar-curvature bound.  Neither kind of
input is inferred from small integral Einstein defect.

Two-forms have their exterior-algebra norm, and $W^+$ has its
Hilbert--Schmidt norm as an endomorphism of $\Lambda^+$.  The scalar
Laplacian is $\Delta=d^*d\geq0$.  In particular, a negative
Kähler--Einstein four-metric with $\Ric=-g$ has scalar curvature $-4$
and self-dual Weyl eigenvalues $(-2/3,1/3,1/3)$.

\subsection{All ordinary classes on one fixed exterior}

We first record the scalar specialization of the monopole identity.
In the normalization $|F_A^+|^2=|\Phi|^4/8$, an unperturbed monopole
solution on a metric of scalar curvature $-4$ satisfies
\[
 0\leq y:=|\Phi|^2\leq4,
 \qquad 4\int|\nabla_A\Phi|^2=\int(4y-y^2).
\]
If the determinant class $K$ has $K^2=1$, Chern--Weil gives
$\int y^2=32\pi^2+8\|F_A^-\|_2^2$.  Consequently
\begin{equation}\label{eq:eqg-sw-defect}
 16(V-2\pi^2)
 =\|4-y\|_2^2+8\|\nabla_A\Phi\|_2^2+8\|F_A^-\|_2^2.
\end{equation}
Indeed, expanding the first square and using the preceding two
identities gives exactly $16V-32\pi^2$.  No Einstein assumption is
needed in this identity.

\begin{theorem}[Identification of every ordinary monopole class]
\label{thm:eqg-ordinary}
Let $N$ be a closed, simply connected, oriented smooth four-manifold.
Suppose that $K^2=1$ and that $K$ and $-K$ are ordinary monopole classes.
Assume that there are metrics $g_j$ with
\[
 s_{g_j}=-4,\qquad V_j\longrightarrow2\pi^2,
\]
and marked smooth convergence on compact subsets of the regular locus
of a fixed negative Kähler--Einstein orbifold $(X,g_0,\omega_0)$
with one Wahl singularity.  Let $C\subset X_{\rm reg}$ be the compact
core obtained by removing a neighborhood of that point, identified
with a submanifold of $N$.  Assume that
\begin{equation}\label{eq:eqg-core-cohomology}
 H^2(N;\mathbb R)\longrightarrow H^2(C;\mathbb R)
 \text{ is an isomorphism},\qquad
 K|_C=[\omega_0/(2\pi)].
\end{equation}
Then the complete set of ordinary monopole classes is $\{K,-K\}$.
\end{theorem}

\begin{proof}
Fix an ordinary monopole class $a$.  On the same metric $g_j$, take
solutions in classes $a$ and $K$.  Choose the simultaneous curvature
sign convention so that $f_{a,j}/(2\pi)$ and $f_{K,j}/(2\pi)$ represent
these real classes.  The maximum principle gives
$|f_{a,j}^+|,|f_{K,j}^+|\leq\sqrt2$.  For this fixed $a$,
\[
 \|f_{a,j}^-\|_2^2
 =\|f_{a,j}^+\|_2^2-4\pi^2a^2
 \leq2V_j-4\pi^2a^2=O_a(1).
\]
This bound remains valid when $a^2<0$.
Equation~\eqref{eq:eqg-sw-defect} gives
$\|f_{K,j}^-\|_2\to0$.  Hence
\[
 4\pi^2a\cdot K
 =\int\langle f_{a,j}^+,f_{K,j}^+\rangle
       -\int\langle f_{a,j}^-,f_{K,j}^-\rangle
\]
implies $|a\cdot K|\leq1$.
Both classes are characteristic.  Since $N$ is simply connected,
$H^2(N;\mathbb Z)$ is torsion-free, and $a-K=2x$ for an integral
class $x$.  Thus $a\cdot K=1+2x\cdot K$ is odd.  After changing the
sign of $a$, assume $a\cdot K=1$.

Put $D_j=2-\langle f_{a,j}^+,f_{K,j}^+\rangle\geq0$.
The preceding pairing is saturated, so $\int D_j\to0$, and
\[
 |f_{a,j}^+-f_{K,j}^+|^2\leq2D_j,
 \qquad D_j\geq(4-|\Phi_{a,j}|^2)/2.
\]
The bounded amplitude therefore satisfies $|\Phi_{a,j}|^2\to4$
in $L^2$, and the spinor Weitzenböck identity gives
$\|\nabla_{A_{a,j}}\Phi_{a,j}\|_2\to0$.

It is essential to identify the full exterior curvature, not just its
self-dual part.  On a fixed contractible regular ball, abelian Coulomb
gauge and the full $L^2$ curvature bound give a $W^{1,2}$ bound for the
connection.  The pointwise spinor bound then gives an ordinary
$W^{1,2}$ bound for the spinor.  After a subsequence the connections
converge strongly in $L^p$ for $p<4$.  Compactness in $L^2$, together
with the uniform pointwise spinor bound, gives strong spinor convergence
in every finite $L^p$.  These convergences pass to the covariant
derivative equation.  Its limit is a nonzero parallel spinor whose
quadratic form has the same direction as that of the limiting
$K$-solution.

On a Kähler--Einstein ball, write the local Spin$^c$ bundle as the
canonical bundle of spinors tensored with a line bundle.  A parallel
spinor lies in the degree-zero or degree-two line.  The parallel-section
equation determines the determinant curvature to be the canonical
curvature or its negative.  The direction of its quadratic form fixes
which sign occurs.  Thus the two limiting determinant curvatures
coincide.  In particular,
\[
 f_{a,j}-f_{K,j}\rightharpoonup0
 \quad\hbox{in }L^2_{\rm loc}(X_{\rm reg}).
\]
The same conclusion holds for every convergent subsequence on a
regular ball.

Pair with a closed two-form supported in $\operatorname{int}C$.
The resulting topological pairing of $(a-K)|_C$ is zero.
Poincaré--Lefschetz duality and
\eqref{eq:eqg-core-cohomology} imply $a=K$ in real cohomology.
Torsion-freeness gives equality of integral classes.  Restoring the
possible sign proves the theorem.
\end{proof}

The core isomorphism is satisfied, for example, when the marked
complement is a rational homology ball with rational-homology-sphere
boundary.  This observation does not construct the required metric
sequence.  A numerical Yamabe invariant, or knowledge of just the
nonzero Seiberg--Witten basic classes, cannot replace the fixed-exterior
and cohomological hypotheses of Theorem~\ref{thm:eqg-ordinary}.

\subsection{Hodge correction without a positive spectral gap}

\begin{lemma}[A closed form and its self-dual harmonic part]
\label{lem:eqg-hodge}
Let $\Omega$ be a closed real two-form on a closed oriented
four-manifold, and let $\alpha$ be its harmonic representative.
Set $\psi=\alpha^+$.  Then
\begin{equation}\label{eq:eqg-hodge}
 \|\psi-\Omega^+\|_2^2+\|\alpha^-\|_2^2
 =\|\Omega^-\|_2^2.
\end{equation}
In particular, each term on the left is bounded by the right, without
any assumption on the first positive two-form eigenvalue.
\end{lemma}

\begin{proof}
Write $\Omega=\alpha+d\xi$.  Stokes gives
$\int d\xi\wedge d\xi=0$, hence
$\|(d\xi)^+\|_2=\|(d\xi)^-\|_2$.
The harmonic form $\alpha^-$ is orthogonal to the exact form $d\xi$,
and thus to $(d\xi)^-$.  Since
$\psi-\Omega^+=-(d\xi)^+$, Pythagoras proves
\eqref{eq:eqg-hodge}.
\end{proof}

\begin{proposition}[The four-term volume identity]
\label{prop:eqg-four-term}
Suppose $s_g=-4$, $K^2=1$, and $\alpha$ is the harmonic representative
of $2\pi K$.  Assume that $\psi=\alpha^+$ is nowhere zero and that
its induced almost-complex structure has canonical class $K$.
Set $u=|\psi|_g/\sqrt2$ and $\widehat g=ug$.  Then
\begin{equation}\label{eq:eqg-four-term}
 V-2\pi^2
 =\|u-1\|_2^2+3\|d\sqrt u\|_2^2
 +\frac14\int|\widehat\nabla\psi|_{\widehat g}^2
                      \,d\mu_{\widehat g}
 +\frac12\|\alpha^-\|_2^2.
\end{equation}
Unmarked norms use $g$.  Only the self-dual part $\psi$, not the
entire class representative $\alpha$, is the almost-Kähler form.
\end{proposition}

\begin{proof}
The form $\psi$ is closed and has $\widehat g$-norm $\sqrt2$.
Write
\[
 U_r=\int u^r\,d\mu_g,\quad
 H=\|d\sqrt u\|_2^2,\quad R=\|\alpha^-\|_2^2,\quad
 T=\int|\widehat\nabla\psi|^2_{\widehat g}\,d\mu_{\widehat g}.
\]
The intersection pairing gives $2U_2-R=4\pi^2$.
The almost-Kähler canonical-connection formula is
\[
 \int\left(s_{\widehat g}
            +\frac12|\widehat\nabla\psi|^2\right)d\mu_{\widehat g}
 =4\pi c_1(TN)\cdot[\psi]=-4U_2.
\]
Here $c_1(TN)=-K$ and
$\int\alpha\wedge\psi=\int\psi\wedge\psi=2U_2$.
The conformal scalar formula gives
$\int s_{\widehat g}\,d\mu_{\widehat g}=6H-4U_1$.
It follows that $U_1-U_2=3H/2+T/8$.
Substitution in $\|u-1\|_2^2=U_2-2U_1+V$ proves
\eqref{eq:eqg-four-term}, including the anti-self-dual term.
\end{proof}

\subsection{A first-derivative estimate with bounded bundle curvature}

\begin{lemma}\label{lem:eqg-bundle-gradient}
Assume the scalar Sobolev inequality
\begin{equation}\label{eq:eqg-sobolev}
 \|f\|_4^2\leq S(\|df\|_2^2+\|f\|_2^2)
\end{equation}
on a closed four-manifold.  Let $E$ be a metric vector bundle with
compatible connection, $|\Ric|\leq K_R$, and $|F^E|\leq K_F$.
If $\nabla^*\nabla v=H$, then
\begin{equation}\label{eq:eqg-bundle-gradient}
 \|\nabla v\|_\infty
 \leq C(S,K_R,K_F)
 \bigl(\|\nabla v\|_2+\|H\|_\infty
                         +K_F\|v\|_\infty\bigr).
\end{equation}
No derivative bound on $H$ or $F^E$, no full-curvature bound, and no
upper bound for total volume is required.
\end{lemma}

\begin{proof}
Put $U=\nabla v$ and use $[\nabla_i,\nabla_j]v=F^E_{ij}v$.
One commutation gives
\[
 \nabla^*\nabla U=\nabla^*T+AU,
 \qquad T_{ij}=-g_{ij}H+F^E_{ij}v,
 \qquad (AU)_j=-\Ric_j{}^kU_k-F^E_{ij}U^i.
\]
Thus $|A|\leq C(K_R+K_F)$ and
$T_0:=\|T\|_\infty\leq C(\|H\|_\infty+K_F\|v\|_\infty)$.
The derivatives of $H$ and $F^E$ remain inside a divergence.

Here is the needed bounded-source estimate.  Write $w=|U|$ and
$w_k=(w-k)_+$.  On $\{w>k\}$ write $U=we$, $|e|=1$,
and test the equation with $w_ke$.  The principal term is
\[
 |dw_k|^2+w_kw|\nabla e|^2.
\]
The angular part of the divergence source is at most
$T_0w_k|\nabla e|$ and is absorbed by half the second term plus
$C T_0^2w_k/w$.  The radial part is treated by Young's inequality.
With $A_k=\{w>k\}$, the Sobolev inequality therefore gives
\[
 \|w_k\|_4^2
 \leq C\left((T_0^2+k^2)|A_k|
                     +\int_{A_k}w_k^2\right).
\]
Choose $k_0=C_0\|U\|_2$ with $C_0$ sufficiently large.
Then $|A_{k_0}|\leq C_0^{-2}$, and Hölder's inequality absorbs the
last term for every $k\geq k_0$.  For $h>k\geq k_0$ this yields
\[
 |A_h|\leq
 \frac{C(T_0^2+k^2)^2}{(h-k)^4}|A_k|^2.
\]
Apply this to $k_n=k_0+(k_0+T_0)(1-2^{-n})$.
Increasing $C_0$ once makes the initial measure small enough for the
quadratic iteration $Y_{n+1}\leq C16^nY_n^2$ to tend to zero.
Thus $w\leq2k_0+T_0$ almost everywhere.  If $U=0$ the conclusion
is immediate.  Smooth approximation of $|U|$ justifies the tests at
its zero set.  This proves \eqref{eq:eqg-bundle-gradient}.
Only a single curvature commutation was used; the final iteration
requires no curvature control for $T^*N\otimes E$.
\end{proof}

\subsection{The exact Einstein endpoint and its first-order energies}

Fix a closed, simply connected oriented $N$ with
$\chi(N)=47$, $\tau(N)=-31$ (hence $b_+(N)=7$), and a monopole
class $K$ such that
\begin{equation}\label{eq:eqg-endpoint-topology}
 K^2=2\chi(N)+3\tau(N)=1.
\end{equation}
Consider an exact sequence
\begin{equation}\label{eq:eqg-endpoint-sequence}
 \Ric(g_j)=-g_j,\qquad
 d_j:=\Vol(g_j)-2\pi^2>0,\qquad d_j\longrightarrow0,
 \qquad\diam(g_j)\leq D,
 \qquad\varepsilon_j:=\sqrt{d_j}.
\end{equation}
The parameter $\varepsilon_j$ is not an Einstein residual: the residual
is identically zero.  These are the topological data of the specified
rational blowdown of a $(-4)$-section in $E(4)$.  The arguments below
use the displayed hypotheses, not a classification of its possible
limits or an assertion that the exact sequence exists.

We use the classical noncollapsed compactness theorem for exact Einstein
four-metrics: a subsequence converges smoothly off finitely many points
to a compact Einstein orbifold, with volume convergence.  Its hypotheses
hold here.  The monopole inequality bounds volume below; Ricci curvature
and diameter then give uniform local noncollapse and
\eqref{eq:eqg-sobolev}; the Einstein Gauss--Bonnet identity gives a fixed
full-curvature $L^2$ bound.  This is the exact Einstein compactness
input, not a smooth compactness assertion for arbitrary bounded-Ricci
metrics; see \href{https://arxiv.org/pdf/0810.4830}{Anderson, Theorem 5.3}.

\begin{proposition}[Global control of one Weyl half]
\label{prop:eqg-half-weyl}
For a sequence \eqref{eq:eqg-endpoint-sequence}, the limit is negative
Kähler--Einstein.  After discarding finitely many terms,
\begin{equation}\label{eq:eqg-half-weyl}
 \|W_j^+\|_\infty\leq C,\qquad
 \|\nabla W_j^+\|_\infty\longrightarrow0,
 \qquad
 \operatorname{spec}(W_j^+)\longrightarrow(-2/3,1/3,1/3)
 \text{ uniformly}.
\end{equation}
No bound for $W_j^-$ is asserted.
\end{proposition}

\begin{proof}
The spinor quadratic forms $\beta_j=-iF_{A_j}^+$ satisfy
\[
 |\beta_j|\leq\sqrt2,\qquad
 \|\nabla\beta_j\|_2^2\leq Cd_j,\qquad
 \bigl\||\beta_j|-\sqrt2\bigr\|_2^2\leq2d_j
\]
by \eqref{eq:eqg-sw-defect} and the differential of the quadratic map.
On each fixed regular core, Rellich compactness and the smooth metric
convergence give a parallel limit of length $\sqrt2$.
On an orbifold covering ball this bounded parallel form extends across
the center: a cutoff at radius $r$ contributes at most $Cr^3$ to its
distributional first derivative.  Elliptic regularity then extends the
parallel equation smoothly.  The resulting parallel orthogonal complex
structure makes the limit Kähler--Einstein.

The self-dual Chern--Weil identity is
\[
 \int_N|W_j^+|^2=2\pi^2-\frac{V_j}{3}
 \longrightarrow\frac{4\pi^2}{3}
 =\int_X|W_0^+|^2.
\]
Smooth convergence on regular cores exhausts this energy.  Thus no
$W^+$ energy concentrates at the finitely many singular points.
The Einstein half-Weyl equation is
\begin{equation}\label{eq:eqg-weyl-pde}
 \nabla^*\nabla W^+=2W^++6(W^+)^2-2|W^+|^2I.
\end{equation}
Kato's inequality gives the scalar inequality
$\Delta|W^+|\leq C(|W^+|+|W^+|^2)$ in the weak sense.
Its small-$L^2$ local estimate is
\[
 \sup_{B_{r/2}}|W^+|
 \leq C(r^{-2}+1)\|W^+\|_{L^2(B_r)}.
\]
To obtain it, test with cutoff powers of $|W^+|$: the quadratic
coefficient is absorbed using its small $L^2$ norm and
\eqref{eq:eqg-sobolev}, first raising the exponent above the critical
one, after which the usual Sobolev iteration applies.  All constants
use the displayed scalar inequality and the uniform local Sobolev
constant, not $W^-$.  The absence of atoms, and smooth convergence
away from the singular set, give a fixed finite covering by such
small-energy balls.  This proves the global $W^+$ bound.

Pairing \eqref{eq:eqg-weyl-pde} with $W^+$ gives
\begin{equation}\label{eq:eqg-weyl-energy}
 \int|\nabla W^+|^2=2\int|W^+|^2+18\int\det W^+\longrightarrow0.
\end{equation}
The cubic integral converges because $W^+$ is now uniformly bounded
and the exceptional neighborhoods have uniformly small volume.
Treat $W^+$ as a section of $\operatorname{Sym}_0(\Lambda^+)$.
The curvature of this bundle depends only on $W^+$ and the fixed
scalar curvature.  The Einstein connection on $\Lambda^+$ is
Yang--Mills.  In the first commutation of
\eqref{eq:eqg-weyl-pde}, its curvature-divergence term therefore
vanishes and the contracted cotangent term is Ricci.  For
$U=\nabla W^+$ one obtains $|\nabla^*\nabla U|\leq C|U|$.
Sobolev iteration gives $\|U\|_\infty\leq C\|U\|_2\to0$.
Finally compare spectra along paths of length at most $D$ with a
regular base point.  Parallel transport preserves eigenvalues, so
this derivative estimate gives uniform spectral convergence.
\end{proof}

The lowest eigenline is now simple.  Simple connectivity makes this
real line bundle orientable.  Choose its global eigenform $\omega_j$
of squared length two, with signs agreeing on a fixed regular core,
and write
\begin{equation}\label{eq:eqg-eigenvalues}
 \lambda_1=-2/3-q,\qquad
 \lambda_2=1/3+q/2-b,\qquad
 \lambda_3=1/3+q/2+b,\qquad b\geq0.
\end{equation}
The associated almost-complex structure will be denoted by $J_j$.

\begin{proposition}[First-order energy ratios]
\label{prop:eqg-ratios}
For \eqref{eq:eqg-endpoint-sequence},
\begin{gather}
 \|q_j\|_\infty\leq Cd_j,\qquad
 \|b_j\|_\infty+\|\nabla W_j^+\|_\infty
                  +\|\nabla\omega_j\|_\infty
 \leq C\sqrt{d_j},                                      \label{eq:eqg-rates}\\
 \frac{\int b_j^2}{d_j}\longrightarrow\frac12,
 \quad
 \frac{\int|\nabla W_j^+|^2}{d_j}\longrightarrow10,
 \quad
 \frac{\int|\nabla\omega_j|^2}{d_j}\longrightarrow4,
 \quad
 \frac{\int|d\omega_j|^2}{d_j}\longrightarrow0.          \label{eq:eqg-ratios}
\end{gather}
\end{proposition}

\begin{proof}
Complete $\omega_1=\omega$ locally to a quaternionic frame
$(\omega_1,\omega_2,\omega_3)$ of squared lengths two, and write
\[
 \nabla\omega_1=\beta\otimes\omega_2+\gamma\otimes\omega_3,
 \quad
 \nabla\omega_2=-\beta\otimes\omega_1+\alpha\otimes\omega_3,
 \quad
 \nabla\omega_3=-\gamma\otimes\omega_1-\alpha\otimes\omega_2.
\]
The curvature of the projected canonical-line connection gives a
global closed form
\begin{equation}\label{eq:eqg-canonical-form}
 \Omega=d\alpha=(1+q)\omega-\beta\wedge\gamma,
 \qquad [\Omega/(2\pi)]=K_J,
 \qquad |\Omega^-|\leq\tfrac14|\nabla\omega|^2.
\end{equation}
This is a local connection formula: it does not say that $\Omega$
is globally exact.  The formula follows by projecting the
$\Lambda^+$ curvature matrix $W^+-I/3$ onto the transverse rotation
line; its lowest entry is $-1-q$.  The last bound follows from
$|\nabla\omega|^2=2(|\beta|^2+|\gamma|^2)$.
By Proposition~\ref{prop:eqg-half-weyl}, $\Omega$ is uniformly close
to $\omega$, hence eventually symplectic and taming $J$.
The canonical square is $K_J^2=2\chi+3\tau=1$.

Let $\alpha_0$ be the harmonic representative of $\Omega$ and
$\psi=\alpha_0^+$.  Lemma~\ref{lem:eqg-hodge} gives
$\|\psi-\omega\|_2\to0$.  The harmonic two-form equation
\[
 \nabla^*\nabla\psi=(2W^++4I/3)\psi
\]
and Sobolev iteration first give a uniform pointwise bound.
Its integrated equation gives $\|\nabla\psi\|_2\to0$,
since $\psi$ approaches the lowest axis in $L^2$.
Commuting once, using the Yang--Mills property as above, gives a
bounded-coefficient equation for $\nabla\psi$ with source
$2(\nabla W^+)\psi$.  Consequently
$\|\nabla\psi\|_\infty\to0$.
A point where $|\psi-\omega|$ is small exists by the $L^2$ bound
and the volume lower bound.  Comparison along paths of length at
most $D$ makes this closeness uniform.  Thus $\psi$ is nowhere zero
and determines an almost-complex structure homotopic to $J$.
Proposition~\ref{prop:eqg-four-term} applies with $K_J$; in particular
$u=|\psi|/\sqrt2\to1$ uniformly and
$\widehat T:=\int|\widehat\nabla\psi|^2\leq4d$.

In the conformal metric $\widehat g=ug$, decompose $W^+$ relative
to the unit axis $\psi/\sqrt2$ as
\[
 W^+_{\widehat g}=
 \begin{pmatrix}a&v^{\mathsf T}\\v&-aI/2+\widehat B\end{pmatrix},
 \qquad\tr\widehat B=0,
 \quad w=W^+_{\widehat g}(\psi,\psi)=2a,
 \quad t=|\widehat\nabla\psi|^2.
\]
The classical weighted harmonic-Weyl identity for an Einstein
conformal class is
\[
 \int\left(8|W^+|^2
       -4|W^+(\psi)^\perp|^2-s_{\widehat g}w\right)
             u^{-1/2}\,d\mu_{\widehat g}=0;
\]
see \href{https://arxiv.org/pdf/1408.1078}{LeBrun, equation (3)}.
Since $w-s_{\widehat g}/3=t/2$, substituting
$|W^+|^2=3a^2/2+2|v|^2+|\widehat B|^2$ and
$|W^+(\psi)^\perp|^2=2|v|^2$ gives
\begin{equation}\label{eq:eqg-weighted-weyl}
 \int(|v|^2+|\widehat B|^2)u^{-1/2}\,d\mu_{\widehat g}
 =-\frac3{16}\int wt\,u^{-1/2}\,d\mu_{\widehat g}.
\end{equation}
The right side is $O(d)$.  The two eigenvalues near $1/3$ split by
at most a constant times the norm of the off-axis block, by the
eigenvalue perturbation inequality.  Since $u$ stays bounded above
and below, this gives $B_2:=\int b^2\leq Cd$.

Put $Q=\int q^2$ and $G=\int|\nabla W^+|^2$.
The topological energy identity and
\eqref{eq:eqg-weyl-energy}, evaluated on
\eqref{eq:eqg-eigenvalues}, give exactly
\begin{equation}\label{eq:eqg-algebraic-energies}
 2\int q=-d-\tfrac32Q-2B_2,
 \qquad
 G=d+\tfrac92\int(1+q)(4b^2-q^2).
\end{equation}
Thus $G\leq Cd$.  The first-commutator estimate in the proof of
Proposition~\ref{prop:eqg-half-weyl} gives
$\|\nabla W^+\|_\infty\leq C\sqrt d$.
The simple-eigenline derivative formula gives the same bound for
$\nabla\omega$.  The eigenvalue functions are Lipschitz with
derivative at most $|\nabla W^+|$; using $B_2\leq Cd$ and diameter
then gives $\|b\|_\infty\leq C\sqrt d$.

Let $\nu=\omega/\sqrt2$.  Differentiating its eigenvalue equation
twice, and substituting \eqref{eq:eqg-weyl-pde}, gives
\begin{equation}\label{eq:eqg-radial-mass}
 (\Delta+2)q=4b^2-3q^2
       -2\sum_{k=2,3}(\lambda_k-\lambda_1)
                              |(\nabla\nu)_k|^2.
\end{equation}
For clarity, the last term is the usual simple-eigenvalue correction:
the second derivative of the eigenvalue is its differentiated matrix
entry minus twice the squared derivative of the eigenvector, weighted
by the other eigenvalue differences.  All entries are smooth even
when the two transverse eigenvalues coincide.  The inverse maximum
principle for $\Delta+2$, the preceding bounds, and $q\to0$ give
$\|q\|_\infty\leq Cd+C\|q\|_\infty^2\leq Cd$ after absorption.
This proves \eqref{eq:eqg-rates}.

Equation~\eqref{eq:eqg-canonical-form} and
Lemma~\ref{lem:eqg-hodge} now imply
\[
 \|\alpha_0^-\|_2+\|u-1\|_2=O(d),
 \qquad\|\psi-\omega\|_2=O(d).
\]
The four-term identity bounds $\|\nabla\psi\|_2$ by $C\sqrt d$;
this follows also by the conformal connection formula and its
$\|d\sqrt u\|_2$ term.  The once-commuted harmonic equation therefore
gives $\|\nabla\psi\|_\infty=O(\sqrt d)$.
The angle between $\psi$ and the lowest axis is $O(d)$ in $L^2$
and $O(\sqrt d)$ uniformly, by diameter comparison.
If $\mu=\langle W^+\psi,\psi\rangle/|\psi|^2$, it follows that
$\mu+2/3=O(d)$.  The norm equation for the harmonic form reads
\[
 2u\Delta u=4u^2(\mu+2/3)-|\nabla\psi|^2+2|du|^2.
\]
It yields $\|\Delta u\|_\infty=O(d)$.  Integration by parts gives
$\int|du|^2=\int(u-1)\Delta u=O(d^2)$.
The four-term identity consequently becomes
$\widehat T=4d+O(d^2)$.

In the block decomposition, the angle estimate gives
$\int|v|^2=O(d^2)$, while
$\int|\widehat B|^2=2B_2+o(d)$.
The latter follows by rotating the nearly degenerate transverse
plane, using its $O(d)$ $L^2$ angle error, and $u\to1$;
all conformal weights also tend uniformly to one.
Since $w\to-4/3$, equation~\eqref{eq:eqg-weighted-weyl} gives
$2B_2=d+o(d)$.  Then \eqref{eq:eqg-algebraic-energies} gives
$G=10d+o(d)$.

Integrating \eqref{eq:eqg-radial-mass} and using
\eqref{eq:eqg-algebraic-energies} yields
\[
 2\int\sum_{k>1}(\lambda_k-\lambda_1)|(\nabla\nu)_k|^2
 =d+6B_2-\tfrac32Q.
\]
The eigenvalue differences tend uniformly to one, so
$\int|\nabla\omega|^2=4d+o(d)$.
Finally the self-dual Hodge--Weitzenböck formula gives the exact identity
\[
 d+\|d\omega\|_2^2+\tfrac32Q+2B_2
 =\tfrac12\|\nabla\omega\|_2^2.
\]
Substitution proves the last limit in \eqref{eq:eqg-ratios}.
\end{proof}

These positive limiting energies are compatible with each other.
They are necessary conditions on a putative exact endpoint sequence,
not a contradiction and not a universal near-Einstein estimate.

\subsection{The holomorphic first-order tensor}

\begin{proposition}\label{prop:eqg-holomorphic}
Under \eqref{eq:eqg-endpoint-sequence}, the fixed orbifold limit $X$
carries a nonzero tensor
\[
 \sigma\in H^0_{\rm orb}(X,\Omega^1\otimes K_X).
\]
If $B_j$ denotes the trace-free transverse block in
\eqref{eq:eqg-eigenvalues}, its normalized regular-locus limit
$B_1=\lim B_j/\varepsilon_j$ corresponds to $\partial_{K_X}\sigma$
under a fixed nonzero real bundle isomorphism.  In particular their
vanishing at any regular or orbifold point is equivalent.
\end{proposition}

\begin{proof}
Set $A_j=\nabla\omega_j/\varepsilon_j$.
Equation~\eqref{eq:eqg-rates} bounds $A_j$ uniformly.  On a fixed
regular core, the once-commuted Einstein equation and the spectral
projection formula give uniform bounds for all derivatives of $A_j$:
the metric coefficients converge smoothly there, and the homogeneous
equation for $\nabla W^+$ has right side linear in $\nabla W^+$.
Local elliptic estimates and differentiation therefore retain the
factor $\varepsilon_j$.  Extract a smooth local limit $A$.
The exceptional neighborhoods have volume $O(r^4)$ uniformly, so
the pointwise bound prevents loss of its normalized energy.  By
\eqref{eq:eqg-ratios},
\[
 \int_X|A|^2=4,\qquad \operatorname{Alt}A=0,
 \qquad A\in T^*X\otimes\omega_0^\perp.
\]
Complexifying, the alternation map
$\Lambda^{0,1}\otimes\Lambda^{2,0}\to\Lambda^{2,1}$ is an
isomorphism in complex dimension two.  Thus
$A=\sigma+\overline\sigma$ with
$\sigma\in\Omega^{1,0}\otimes K_X$.

For $J_j$-type vectors $Z$ and $\overline U$, the exact Einstein
curvature identity and the lowest-eigenline condition give
\[
 R^{\Lambda^+}_{g_j}(Z,\overline U)\omega_j=0.
\]
Indeed the self-dual part of a $(1,1)$ bivector is in the lowest-axis
line, and the block-diagonal curvature operator preserves that line;
its infinitesimal rotation fixes $\omega_j$.
Divide $d_\nabla\nabla\omega_j=R^{\Lambda^+}\omega_j$ by
$\varepsilon_j$ and take its mixed-type part.  Its $K_X$ component
is $\nabla^{0,1}\sigma=0$.  The tensor is bounded near each
orbifold point.  On a smooth covering ball its holomorphic
coefficients extend by Hartogs, equivariantly.  The positive energy
above proves that $\sigma\ne0$.

The $(2,0)$ part of the same curvature equation is the transverse
matrix $B_j+q_jI/2$.  The scalar part disappears after division by
$\varepsilon_j$, since $q_j=O(\varepsilon_j^2)$.
In a unitary transverse frame, a trace-free symmetric matrix has
entries $\left(\begin{smallmatrix}x&y\\y&-x\end{smallmatrix}\right)$;
its action on the complex transverse line is represented by the
nonzero complex coefficient $x+iy$, up to the fixed frame convention.
The exterior covariant derivative on the limiting side is
$\partial_{K_X}\sigma$.  This gives the asserted real bundle
isomorphism and the vanishing equivalence, without a phase-dependent
identification of their numerical coefficients.
\end{proof}

For a merely approximate Einstein sequence the mixed curvature
equation has a trace-free Ricci source.  Division by
$\varepsilon_j$ need not make it vanish.  Thus the holomorphic
conclusion above cannot be transferred without a relative error
estimate and the stated compactness hypotheses.

\subsection{Critical norms on a finite bubble tree}

\begin{lemma}[The physical-scale Hessian sum]
\label{lem:eqg-tree}
Suppose a sequence has a fixed finite noncollapsed Einstein bubble-tree
decomposition.  Assume that the rescaled charts on each core and
dyadic neck block have uniform smooth geometric bounds, and that a
tensor $J_j$ obeys
\[
 |\nabla^2J_j|\leq C\varepsilon_j/r
\]
on a block of physical scale $r$.  Core volumes are at most $Cr^4$,
and each neck has the usual dyadic scales between the physical
parent and child scales.  Then
\begin{equation}\label{eq:eqg-tree}
 \|\nabla^2J_j\|_{L^p}\leq C_p\varepsilon_j\quad(1\leq p<4),
 \qquad
 \|\nabla^2J_j\|_{L^{4,\infty}}\leq C\varepsilon_j.
\end{equation}
After fixing the outer scale, the corresponding strong estimate is
\[
 \|\nabla^2J_j\|_4^4
 \leq C\varepsilon_j^4
 \left(1+\sum_{u\to v}\log\frac{a_{u,j}}{a_{v,j}}\right),
\]
where $a_{v,j}$ are cumulative physical node scales.
\end{lemma}

\begin{proof}
A core of scale $a$ contributes at most
$C_p\varepsilon_j^pa^{4-p}$ to the $p$th power of the norm.
A neck contributes at most
$C_p\varepsilon_j^p\sum_k r_k^{4-p}$.
For $p<4$ the dyadic geometric series is bounded by its largest
scale, uniformly in the number of annuli.  There are only finitely
many cores and necks.  This proves the first assertion.
For a threshold $t>0$, only blocks with $r\leq C\varepsilon_j/t$
can contribute to $\{|\nabla^2J_j|>t\}$.
On each neck their volumes sum to at most
$C(\varepsilon_j/t)^4$.  A core either contributes nothing or obeys
the same bound.  Summing over the fixed finite tree proves the weak
$L^4$ assertion.  At $p=4$ each neck annulus costs
$C\varepsilon_j^4$, and their number is bounded by a constant plus
the logarithm of the physical scale ratio.  This proves the last
estimate.
\end{proof}

In the endpoint setting, the displayed block bound follows from the
uniform first-derivative bound and the local homogeneous elliptic
equation on these specified rescaled charts.  The lemma does not
supply the charts or a bubble decomposition.  It gives neither a
strong $L^{4,1}$ estimate nor an example attaining the logarithmic
upper bound.

\subsection{A fixed-background critical deformation estimate}

Here the background is a fixed compact smooth complex orbifold of
complex dimension two.  All constants in this subsection belong to
that background.  Let $D=\bar\partial$ on $T^{1,0}$-valued forms,
let $H_q$ and $G_q$ be its harmonic projection and Green operator,
and set
\[
 \mathcal B=W^{1,(4,1)}(\Lambda^{0,1}\otimes T^{1,0}),\qquad
 \mathcal Y=L^{4,1}(\Lambda^{0,2}\otimes T^{1,0}),
 \qquad F(\mu)=D\mu+\tfrac12[\mu,\mu].
\]
The endpoint embedding $\mathcal B\hookrightarrow C^0$, fixed-background
elliptic estimates, and the first-order bracket give
\[
 B:=D^*G_2:\mathcal Y\longrightarrow\mathcal B,
 \qquad
 \|[u,v]\|_{\mathcal Y}\leq C\|u\|_{\mathcal B}\|v\|_{\mathcal B}.
\]
These are strong Lorentz spaces; the weak $L^4$ conclusion of
Lemma~\ref{lem:eqg-tree} is not the required input.

\begin{proposition}[Conditional projection and integrable retraction]
\label{prop:eqg-kuranishi}
Let $f:H^1\to H^2$ be the local Kuranishi obstruction map.
Suppose a fixed linear functional $\ell$ satisfies
\[
 \ell f(h)=a(h)f(h),\qquad a(0)=0,
\]
where $a(h)$ is an analytic family of linear functionals.
For sufficiently small $\mu\in\mathcal B$ with $D^*\mu=0$,
\begin{equation}\label{eq:eqg-kuranishi-projection}
 |\ell H_2F(\mu)|\leq C\|\mu\|_{\mathcal B}\|F(\mu)\|_{\mathcal Y}.
\end{equation}
If, more strongly, every actual obstruction for the smooth Dolbeault
deformation problem vanishes on every Artin small extension, then
arbitrary sufficiently small $\mu\in\mathcal B$ admits $\nu$ with
\begin{equation}\label{eq:eqg-kuranishi-retraction}
 F(\nu)=0,\qquad
 \|\nu\|_{\mathcal B}\leq C\|\mu\|_{\mathcal B},\qquad
 \|\mu-\nu\|_{\mathcal B}\leq C\|F(\mu)\|_{\mathcal Y},
 \qquad D^*(\mu-\nu)=0.
\end{equation}
Smooth $\mu$ gives smooth $\nu$, without a uniform assertion about
all its higher derivatives.
\end{proposition}

\begin{proof}
The map $\mathcal K(\mu)=\mu+B[\mu,\mu]/2$ is analytic with
derivative the identity at zero.  The Banach inverse function theorem
makes it locally analytically invertible with a fixed Lipschitz inverse.
Complex dimension two gives $DB=1-H_2$.
If $D^*\mu=0$, Hodge decomposition yields
\[
 \mathcal K(\mu)=H_1\mu+BF(\mu).
\]
Set $h=H_1\mu$ and $\nu=\mathcal K^{-1}(h)$.
Then $\|\mu-\nu\|_{\mathcal B}\leq C\|F(\mu)\|_{\mathcal Y}$,
and $F(\nu)$ is harmonic, equal to $f(h)$.
Since $H_2D=0$, the bracket difference estimate gives
\[
 \|H_2F(\mu)-f(h)\|
 \leq C\|\mu\|_{\mathcal B}\|F(\mu)\|_{\mathcal Y}.
\]
Also $\|f(h)\|\leq C\|F(\mu)\|_{\mathcal Y}$.
Using $a(h)=O(|h|)$ proves
\eqref{eq:eqg-kuranishi-projection}.

For the stronger assertion let
$P_Z=H_1+DG_0D^*$, the projection to $\ker D$ in degree one.
Then $\mathcal K(\mu)=P_Z\mu+BF(\mu)$.
For a smooth closed $v$, the function
$F(\mathcal K^{-1}(tv))$ is harmonic-valued and analytic in $t$.
If its coefficients of orders less than $n$ vanish, its coefficient
of order $n$ is precisely the harmonic representative of the
obstruction to lifting the preceding formal solution over the
corresponding Artin small extension.  The assumption makes this
coefficient zero.  Induction and analyticity make the entire function
zero near $t=0$.  Smooth closed forms are dense in the closed
subspace of $\mathcal B$, so this remains true there by continuity.
Take $\nu=\mathcal K^{-1}(P_Z\mu)$.
The Lipschitz bound proves both norm inequalities in
\eqref{eq:eqg-kuranishi-retraction}, and
$D^*B=0$, $D^*P_Z=D^*$ give its last assertion.

For a smooth input the equation
$\nu+B[\nu,\nu]/2=P_Z\mu$ first improves regularity to
$W^{1,p}$ for some $p>4$: the continuous coefficient $\nu$ is
small, so the first-order principal perturbation is absorbed in
the fixed-background elliptic $L^p$ estimate.  One can obtain this
estimate by Neumann inversion with that coefficient frozen and
then use uniqueness in $\mathcal B$.  The system
$(D\nu+[\nu,\nu]/2,D^*\nu)=(0,D^*\mu)$ is now a uniformly
elliptic quasilinear system with Hölder coefficients.  Schauder
iteration gives smoothness.
\end{proof}

The obstruction hypothesis here is about the smooth Dolbeault
deformation problem on this orbifold.  Formal local triviality of a
selected quotient-singularity smoothing parameter is not that
hypothesis, and it does not identify two different deformation
complexes.

\subsection{A low-power Dirac gap with its spectral hypotheses retained}

\begin{lemma}\label{lem:eqg-dirac}
In the exact endpoint setting, use the $U(2)$ projection
$\nabla^U=\nabla-\tfrac12J\nabla J$ and its induced determinant
connection on
\[
 S_m=\Lambda^{0,*}\otimes K_J^m,
 \qquad L_m=K_J^{2m-1},\qquad m=1,2.
\]
Suppose the regular-locus limit is the corresponding canonical
Dirac operator on $X$, and
$H^{0,1}_{\rm orb}(X,K_X^m)=0$.
Then, along this fixed convergent subsequence, there is $\lambda_m>0$
such that
\[
 \|\mathcal D^-_{m,j}\eta\|_2^2
 \geq\lambda_m\|\eta\|_2^2.
\]
When $\tau(N)=-31$, the even kernel dimensions are respectively
$4$ and $5$.
\end{lemma}

\begin{proof}
For $P_{10}=(I-iJ)/2$, the curvature of the projected canonical
connection is
\[
 F^{K_J}=-\operatorname{Tr}\left(
 P_{10}R^{TN}P_{10}
 +P_{10}(\nabla P_{10}\wedge\nabla P_{10})P_{10}\right).
\]
This formula has no independent $\nabla^2J$ term.  The Einstein
lowest-axis condition and \eqref{eq:eqg-rates} give, in the consistent
canonical sign convention, $iF^{K_J}=\omega+O(d_j)$.
In particular the zeroth-order potential in the Dirac Weitzenböck
formula is uniformly bounded for the two fixed values of $m$.
This does not bound the entire negative-spinor connection curvature.

If the asserted gap failed, choose unit eigenspinors for
$(\mathcal D^-_{m,j})^*\mathcal D^-_{m,j}$ with eigenvalues tending
to zero.  The bounded Weitzenböck potential and the uniform Sobolev
inequality give a uniform pointwise bound by scalar Moser iteration.
They also give a uniform covariant $W^{1,2}$ bound.  Smooth local
elliptic convergence on the regular locus yields a limiting odd
harmonic spinor.  The uniform pointwise bound and the $O(r^4)$
volume of exceptional neighborhoods prevent loss of its unit mass.
On an orbifold covering ball, a cutoff at radius $r$ has an error
of $L^2$ norm $O(r)$ in the Dirac equation; removal of the point and
elliptic regularity therefore extend the limit across the center.
The nonzero limit belongs to $H^{0,1}_{\rm orb}(X,K_X^m)$,
contradicting the hypothesis.

The gap eliminates the odd kernel.  The Spin$^c$ index is
\[
 \operatorname{ind}\mathcal D^+_m
 =\frac{(2m-1)^2K_J^2-\tau(N)}8.
\]
For $K_J^2=1$ and $\tau=-31$ this equals
$4+m(m-1)/2$, giving $4$ and $5$.
\end{proof}

These are limit-dependent spectral gaps.  The dimension four at
$m=1$ is a spinor-kernel dimension, not a claim that there are four,
or three, independent holomorphic canonical sections before separating
the degree-two component.  No quadratic improvement of a multiplication
error follows from this lemma.

\subsection{A bubble-period estimate retaining the radial term}

\begin{proposition}\label{prop:eqg-bubble-period}
Let $g_j$ be exact Einstein metrics with $\Ric(g_j)=-g_j$.
On a fixed compact core, possibly after finite covering, assume
\[
 a_j^{-2}\Phi_j^*g_j\longrightarrow k\quad\hbox{smoothly},
 \qquad a_j\longrightarrow0,
\]
where the core lies in a nonflat hyperKähler ALE space and has
$H^1=0$.  Suppose a simple lowest self-dual Weyl axis is represented
by $\omega_j$, $|\omega_j|^2=2$, and write
\[
 W_j^+=(-2/3-q_j)P_j+(1/3+q_j/2)(I-P_j)+B_j,
 \qquad B_j|_{P_j^\perp}\in\operatorname{Sym}_0(P_j^\perp).
\]
All the following bounds are in the physical metric on the chosen
core:
\[
 Q_j=\|q_j\|_\infty\to0,\quad
 M_j=\|B_j\|_\infty,\quad
 \delta_j=\|\nabla\omega_j\|_\infty\to0,\quad
 G_j=\|\nabla B_j\|_\infty+\|dq_j\|_\infty,
 \quad K_j=1+Q_j+\delta_j^2.
\]
Assume a closed two-cycle has nonzero limiting transverse period
vector.  On a fixed smaller core,
\begin{equation}\label{eq:eqg-bubble-period}
 \left||B_j|_{\rm op}-\frac{|q_j|}{2}\right|
 \leq C\left(a_jG_j+a_j\delta_jK_j+a_j^2K_jM_j\right).
\end{equation}
In particular, if $M_j+G_j+\delta_j=O(\varepsilon_j)$ and
$Q_j=O(\varepsilon_j^2)$ with $\varepsilon_j\to0$, then
\begin{equation}\label{eq:eqg-bubble-vanishing}
 \sup|B_j|\leq C(a_j\varepsilon_j+\varepsilon_j^2)
 =o(\varepsilon_j).
\end{equation}
\end{proposition}

\begin{proof}
Use the unit-scale pullback metrics.  The lowest axis tends to a
parallel axis.  Its perpendicular plane bundle is homotopic to the
constant transverse plane of the hyperKähler frame and is therefore
trivial.  In a global oriented self-dual frame write its rotation
connection as $\alpha_1$, with remaining connection components
$\alpha_2,\alpha_3$.  Then
$|\alpha_2|+|\alpha_3|\leq Ca_j\delta_j$.
Choose the compatible curvature signs so that the exact equations are
\begin{align}
 d\alpha_1+\alpha_2\wedge\alpha_3
     &=a_j^2(1+q_j)\omega_1,\label{eq:eqg-bubble-curvature}\\
 \binom{d\alpha_2+\alpha_3\wedge\alpha_1}
       {d\alpha_3+\alpha_1\wedge\alpha_2}
     &=-a_j^2(B_j+q_jI/2)\binom{\omega_2}{\omega_3}.
       \nonumber
\end{align}
In particular $|d\alpha_1|\leq Ca_j^2K_j$.
On this fixed domain impose an abelian Coulomb gauge with normal
boundary condition.  Its harmonic one-form space is zero because
$H^1=0$.  The Hodge estimate in $W^{1,p}$, $p>4$, gives
$\|\alpha_1\|_\infty\leq Ca_j^2K_j$.
The curvature to which this abelian estimate is applied is the closed
form $d\alpha_1$, not the generally nonclosed first component of the
nonabelian curvature.

Integrating the first equation on a cycle, dividing by $a_j^2$,
and passing to the limit shows that its lowest-axis period is zero.
Let $\zeta_j$ be the vector of its two transverse periods;
by hypothesis $|\zeta_j|\geq c>0$.
Stokes applied to the other two equations gives, at any fixed
reference point $x_*$,
\[
 |(B_j(x_*)+q_j(x_*)I/2)\zeta_j|
 \leq C(a_jG_j+a_j\delta_jK_j+a_j^2K_jM_j).
\]
The first term bounds variation of the physical coefficients along
the unit core.  The mixed connection integrals give the second.
The transverse-frame rotation of $B_j$ is its commutator with
$\alpha_1$ and gives the last term; the scalar matrix commutes with
that rotation.  The same estimates transfer from $x_*$ to every
point of the smaller core.

A real trace-free symmetric $2\times2$ matrix satisfies
$B^2=|B|_{\rm op}^2I$.  Therefore
\[
 |(B+tI)\zeta|\geq
 \bigl||B|_{\rm op}-|t|\bigr|\,|\zeta|.
\]
This proves \eqref{eq:eqg-bubble-period}.  Substituting the stipulated
rates and adding $|q_j|/2$ proves
\eqref{eq:eqg-bubble-vanishing}.
\end{proof}

The radial term cannot be omitted.  The matrix
$B=\operatorname{diag}(-q/2,q/2)$ annihilates $(1,0)^{\mathsf T}$
after addition of $qI/2$, even though $B$ need not vanish.  Thus
$q=o(\varepsilon)$, not merely $q=O(\varepsilon)$, is needed for
the stated first-order conclusion from one transverse period.
Two independent transverse periods would give a stronger conclusion,
but $b_2\geq2$ alone does not give their independence.

\begin{corollary}[Vanishing at every actual singular point]
\label{cor:eqg-all-points}
For \eqref{eq:eqg-endpoint-sequence}, the tensor of
Proposition~\ref{prop:eqg-holomorphic} satisfies
\[
 (\partial_{K_X}\sigma)(p)=0
 \quad\hbox{at every orbifold singularity }p\in X.
\]
\end{corollary}

\begin{proof}
Fix a singular point and isolate it by an exterior annulus where
convergence is smooth.  Choose a maximum-curvature point inside.
The maximum tends to infinity; otherwise noncollapse and Einstein
regularity would remove the singularity.  Rescale by the inverse
square root of that maximum.  Its distance from the exterior boundary
tends to infinity in the rescaled metric.  On every fixed rescaled
ball curvature is at most one, and at the base point it is one.
Noncollapse and the exact equation give a complete smooth nonflat
Ricci-flat limit.  The fixed total curvature bound and Euclidean
volume growth give ALE asymptotics.  The uniform $W^+$ bound makes
this limit anti-self-dual.

The classification of Ricci-flat anti-self-dual ALE spaces supplies
a finite simply connected hyperKähler cover.  A nonflat such cover
has a nonzero period of its hyperKähler triple; this is the
nondegeneracy condition in the hyperKähler ALE period classification.
These are the classical ALE inputs, including the quotient case;
see \href{https://arxiv.org/pdf/1102.2442}{Wright, Theorem A and
Sections 2--3}.  The first equation of
\eqref{eq:eqg-bubble-curvature} makes every limiting lowest-axis
period zero.  Some transverse period is therefore nonzero.
Choose a fixed simply connected compact core containing its cycle.

The bounds \eqref{eq:eqg-rates} give
$M_j+G_j+\delta_j=O(\varepsilon_j)$ and
$Q_j=O(\varepsilon_j^2)$ on this core.  The derivative bound for
$B_j$ follows by differentiating its spectral projection; no gap
between its two transverse eigenvalues is required.
Proposition~\ref{prop:eqg-bubble-period} shows that
$|B_j|/\varepsilon_j\to0$ on the core.
Globally $|\nabla(B_j/\varepsilon_j)|\leq C$.
Compare its scalar norm to a regular exterior point at distance
at most $Cr$ from the singular point, first let $j\to\infty$,
then $r\downarrow0$.  This gives $B_1(p)=0$.
Proposition~\ref{prop:eqg-holomorphic} identifies its vanishing with
that of $\partial_{K_X}\sigma$.
Only scalar norms were transported; no synchronization of phases
across the neck, no smoothness of a root bubble, and no comparison
between the volume defect and a bubble scale was used.
\end{proof}

\begin{corollary}[The remaining evaluation hypothesis]
\label{cor:eqg-conditional-gap}
Assume \eqref{eq:eqg-endpoint-topology}, and assume that $N$ admits
no Kähler structure compatible with the chosen orientation.
Suppose that every orbifold limit of every exact sequence
\eqref{eq:eqg-endpoint-sequence} has a set of actual singular points
$S_X$ for which
\[
 H^0_{\rm orb}(X,\Omega^1\otimes K_X)
 \longrightarrow\bigoplus_{p\in S_X}(K_X^2)_p,
 \qquad
 \sigma\longmapsto\bigl((\partial_{K_X}\sigma)(p)\bigr)_{p\in S_X}
\]
is injective.  For every finite $D$ there is $\eta_N(D)>0$ such that
\[
 \Ric(g)=-g,\quad\diam(g)\leq D
 \quad\Longrightarrow\quad
 \Vol(g)\geq2\pi^2+\eta_N(D).
\]
\end{corollary}

\begin{proof}
The scalar monopole inequality gives $V\geq2\pi^2$.
Equality in \eqref{eq:eqg-sw-defect} makes the spinor quadratic form
parallel and nonzero on $N$, contradicting the assumed absence of
a compatible Kähler structure.  Hence every such volume is strictly
larger than $2\pi^2$.
If no positive uniform gap existed, choose metrics with
$0<V_j-2\pi^2<1/j$ and the same diameter bound.  The preceding
propositions produce a nonzero holomorphic tensor on a subsequential
orbifold limit, whereas Corollary~\ref{cor:eqg-all-points} puts it in
the kernel of the asserted injective evaluation map.  This is a
contradiction.  Injectivity on that one extracted limit is sufficient;
no uniform inverse norm over a moduli space is used.
\end{proof}

The evaluation and classification of all possible limits are additional
hypotheses in this final corollary.  They have not been deduced from
the numerical data $K^2=1$, $\chi=47$, $\tau=-31$.
Even under that hypothesis the conclusion concerns exact Einstein
metrics near one volume endpoint with bounded diameter.  It neither
excludes exact Einstein metrics at larger volume nor supplies a
diameter bound for arbitrary saturating sequences.
